\scanpage{34}
% Pencil archive foliation: 33.
mais on fera attention que l'inclusion
\[
\unsure{\mathbb{Q}|\mathbb{Z}}(-1)\otimes \mathcal{M}^{+i}(X)
\hookrightarrow \mathcal{M}^{+(i+2)}(X)
\]
plus gén.
\[
\unsure{\mathbb{Q}|\mathbb{Z}}(-j)\otimes \mathcal{M}^{+i}(X)
\subset \mathcal{M}^{+(i+2j)}(X)
\]
est \uline{stricte} en général, i.e. n'est pas
une équivalence de catégories.

Noter que $\mathcal{M}^{+i}(X)$ est \uline{épaisse} dans $\mathcal{M}^{+j}(X)$
\margin{indices en bas}
\deleted{[elle est définie dans la catégorie}, et de
\unsure{$=$|$\overline{\mathrm m}$} $\mathcal{M}^i(X)$ épaisse dans $\mathcal{M}^j(X)$ $(i\leq j)$.
Les quotients
\[
\inserted{C_i^+(X)=}
\operatorname{Gr}_i^{\unsure{F|\varepsilon}}(\mathcal{M}^+(X))
\deleted{\unsure{\mathord{?}}}
\inserted{\simeq \mathcal{M}_i^+(X)/\mathcal{M}_{i-1}^+(X)}
\quad\text{et}\quad
C_i(X)=\operatorname{Gr}_i^{\unsure{F|\varepsilon}}(\mathcal{M}(X))
\simeq \mathcal{M}_i(X)/\mathcal{M}_{i-1}(X)
\]
sont fort intéressants.
Noter que les $\unsure{C_i(X)|C_{2i}(X)}$ sont tous
\deleted{isomorphes}\inserted{équivalents}
par twisting
\[
C_i(X)\overset{\mathbb{Z}(-j)}{\simeq}C_{i+2j}(X),
\]
en particulier tous \deleted{isomorphes}\inserted{équiv. can.} à $C_0(X)$.
D'ailleurs
\[
C_i^+(X)\underset{\sim}{\subset}C_i(X)\simeq C_0(X)
\]
\inserted{équiv. à une} sous-catégorie pleine et \uline{épaisse} de
$C_i(X)\simeq C_0(X)$.
